% v13 — reconstructed from Eva_JorgeNike_versao12.pdf and updated to match
% the current model implementation (branch evam-model-review-fixes @ b390871):
% empirical Eurostat income anchor, hard income-budget gate, Wright's-law
% price learning, trait rescaling, Portugal 2010-2020 calibration with
% 2021-2024 hold-out, and fully regenerated multi-seed figures
% (outputs/v13_figures/, produced by scripts/make_v13_figures.py).
% Sections 1-2 and all still-valid prose are kept verbatim from v12.
\documentclass[11pt]{article}
\usepackage[utf8]{inputenc}
\usepackage[T1]{fontenc}
\usepackage{amsmath,amssymb}
\usepackage{graphicx}
\usepackage{booktabs}
\usepackage[margin=1.1in]{geometry}
\usepackage[hidelinks]{hyperref}

\title{An agent-based model for electric vehicles adoption with
pro-environmental behaviour dynamics}
\author{The Author}
\date{July 23, 2026 (v13)}

\begin{document}
\maketitle

\section{Introduction}

In this paper, we intent to design an exploratory model to study human
behavior change towards sustainability issue in socio-economics systems. The
aim of our work is to develop an agent-based model (ABM) for electric
vehicles (EVs) adoption using dynamic computer simulation modelling.

The ABMs provide useful tools to take advantage of a more realist description
about human decision-making process (\cite{bonabeau2002}), in particular for
new energy technologies adoption (\cite{rai2016}). We are focused on EV
adoption because this vehicle technology has the greatest potential to
achieve social benefits regarding climate change, increased energy
independence and reduced greenhouse gas (GHG) emissions. The main modelling
challenge is to represent sustainable consumption behaviour (SCB) that give
rise to EV adoption (\cite{nielsen2024}) for meet goals from three
dimensions: economy, ecology and social, known as the triple bottom line
(\cite{elkington1997}).

In many circumstances, consumer decisions are contextual and affected by
psychological processes (or cognitive biases) such as habit, social norms,
bounded rationality, loss aversion, cognitive depletion, temporal
constraints, anchoring, and peer influence (\cite{orourke2015}).

Many studies in the SCB literature (\cite{syed2024}) highlights a persistent
gap between consumer attitudes, intention and behaviour, known as the
``intention-behavior gap''. This barrier prevents widespread EV adoption so
that it requires a multi-dimensional perspective including psychological,
social, economic, environmental, and contextual factors to better capture
the complexity of EV adoption on the ongoing sustainability transition.
% TODO(v12 carry-over): two unresolved citations [?],[?] in this sentence.

\section{Literature Review}
\subsection{Framework for EV adoption models}
% TODO: section body pending (unchanged from v12).

\section{Model Overview}

\subsection{Model Design}

The Electric Vehicle Adoption Model (EVAM) is implemented as an extension of
the Affordance Landscape Model (from now called it base model) by
introducing the electric vehicle (EV) adoption dynamics with charging
infrastructure evolution, and charging access adjusted to congestion. This
model is embedded within a broader behavioural and spatial framework,
preserving the modularity and generality of the original architecture based
on affordances.

While EVAM inherits the spatial grid, scheduling logic, and data-collection
structure from the base model, all mechanisms related with EV are
encapsulated within the extension. These include a population of
heterogeneous EV consumer agents, the computation of total cost of ownership
via an external module (\texttt{ev\_costs.py} file), and the evolution of
public charging infrastructure. This modular design preserves conceptual
independence while enabling focused analysis of EV adoption dynamics.

\subsection{Model Structure}

EVAM is a dynamic model structured in five blocks:

\begin{enumerate}
\item \textbf{Agents.}
  \begin{itemize}
  \item Agents are consumers (households) that will take a decision about EV
    choice when own vehicle has reached or exceeded its replacement
    interval.
  \item Agents show their behavioural state towards environment,
    pro-environmental or non-environmental, which is inherited from
    pro-environmental learning dynamics of the base model.
  \item Agents are endowed with socio-economic attributes (income, mileage,
    vehicle age), behavioural attributes (environmental concern, price
    sensitivity, peer sensitivity, range anxiety), vehicle replacement
    behaviour, and access to home charging.
  \item Agents are randomly located on the spatial grid, interact with their
    local environment and neighbours, and can be connected to a (optional)
    social network.
  \end{itemize}

\item \textbf{Environment (Space).}
  The environment is represented by a two-dimensional discrete grid with
  three spatial layers:
  \begin{enumerate}
  \item[(i)] An affordance landscape layer inherited from the base model,
    representing pro-environmental cues;
  \item[(ii)] A charging access layer representing the availability of
    public charging infrastructure;
  \item[(iii)] A effective charging access layer in which charging access is
    discounted by congestion in order to meet local supply and demand. It
    coincides with charging access layer when charger capacity is infinite.
  \end{enumerate}

  Charging access layer is a spatial field defined on the grid whose purpose
  is to measure accessibility to public charging infrastructure for each
  cell. Public charging infrastructure is dynamic and is characterized by
  exogenous or demand-driven infrastructure expansion. Chargers are placed
  uniformly at random on the grid (in the exogenous mode) or placed in
  high-demand areas (in the demand-driven mode), subject to the constraint
  that each cell may contain at most one charger.

  The growing demand for EVs underscores the critical need for efficient and
  accessible charging infrastructure (\cite{mohammed2024}).

\item \textbf{Behavioral economic analysis related with EVs.}

  Agents start by estimating the total cost of ownership (TCO) for both EV
  and internal combustion engine (ICE) vehicles. TCO includes purchase price
  (without financing costs), operational costs (fuel and electricity prices,
  annual mileage, fuel and energy consumptions), and maintenance costs.

  In our analysis instead of the EV purchase price $p_{EV}$ it is considered
  the effective EV purchase price at time $t$. Effective EV purchase price
  describes cost reduction of technologies with the increase in production
  volume over time, represented by a learning (or experience) curve. The
  model implements Wright's law keyed on cumulative adopters $N_t$:
  \begin{equation}
    p^{e\!f\!f}_{EV,t} \;=\;
    \max\!\Big( p_{EV}\,\big(N_t / N_0\big)^{-b},\;
                p_{EV}\cdot \mathit{floor\_share} \Big),
    \qquad b = -\log_2\!\big(1 - \mathit{LR}\big),
  \end{equation}
  where $\mathit{LR}$ is the fractional cost decline per doubling of
  cumulative adopters (default $\mathit{LR}=0.18$, calibrated;
  cf.\ commonly cited lithium-ion pack learning rates), $N_0$ is the
  reference adopter count at which the list price applies, and
  $\mathit{floor\_share}$ is the minimum EV price fraction. A simpler linear
  proxy $p^{e\!f\!f}_{EV,t} = \max\big(p_{EV}(1-\eta\, s_{EV,t}),\,
  p_{EV}\cdot\mathit{floor\_share}\big)$, with $s_{EV,t}$ the EV adoption
  share and $\eta$ a learning coefficient, is retained as an option for
  backward comparability.

  Let $p_{ICE}$ and $p_{EV}$ purchase prices, $S$ a subsidy for EV, $p_f$
  and $p_e$ fuel and electricity prices respectively, $c_{ICE,i}$ and
  $c_{EV,i}$ fuel and electricity consumptions in $L/km$ and $kWh/km$
  respectively, $m_i$ the mileage in km/year, $Y$ the ownership horizon in
  years, and $\mathit{mc}_{ICE}$ and $\mathit{mc}_{EV}$ annual maintenance
  costs. The $i$ agent's TCO for EV and ICE vehicles respectively, are
  given by
  \begin{align}
    TCO_{EV,i,t}  &= \big(p^{e\!f\!f}_{EV,t} - S_t\big)^{+}
      + a(Y, r)\,\big(c_{EV}\, m_i\, p_{e,t} + \mathit{mc}_{EV}\big),\\
    TCO_{ICE,i,t} &= p_{ICE}
      + a(Y, r)\,\big(c_{ICE}\, m_i\, p_{f,t} + \mathit{mc}_{ICE}\big),
  \end{align}
  where $a(Y,r)=\sum_{k=1}^{Y}(1+r)^{-k}$ is an annuity factor with
  discount rate $r$; the default $r=0$ gives $a(Y,0)=Y$ (undiscounted).
  The subsidy and the fuel and electricity prices may follow exogenous
  per-step schedules (Section~\ref{sec:results} uses the historical
  Portuguese subsidy ramp), so they carry a time index.

  TCO is used to verify cost-competitiveness for each pair of vehicle
  choices, taking ICE vehicle cost as ``reference price''. The TCO gap
  \begin{equation}
    \Delta(TCO_{i,t}) = TCO_{ICE,i} - TCO_{EV,i,t}
  \end{equation}
  allows then to analyse cost-competitiveness because EV are cost-competitive
  if $\Delta(TCO_{i,t})>0$.

\item \textbf{Updating mechanisms.}
  In each step time $t$ the procedure involved agents is the following:
  \begin{enumerate}
  \item[(a)] Agents might change their generic behavioral state
    (pro-environmental or non-environmental) through social learning
    process. It means that agents perform an affordance behaviour loop
    (move, act, and learn);
  \item[(b)] Once fixed vehicle replacement interval $[6,14]$ in years, it
    is incremented vehicle age by one time step until reaches the vehicle
    replacement event;
  \item[(c)] The model resolves the current subsidy, fuel price, and
    electricity price from their (optional) exogenous schedules, and
    updates the effective EV purchase price by the learning curve;
  \item[(d)] The peer adoption share is determined representing a fraction
    of visible peers who currently own an electric vehicle;
  \item[(e)] If social diffusion option is enabled, agents revise both the
    generic pro-environmental state and the specific (regarding EVs)
    environmental concern. This mechanism allows perceptions relevant to EV
    adoption to spread through the social network. The agent's specific
    environmental concern increases proportionally to peer adoption since
    that exposure to EV adopters increases pro-environmental attitudes;
  \item[(f)] Agents reduce perceived range anxiety over time due to peer EV
    adoption.
  \end{enumerate}

  Public charging infrastructure is developed over time through exogenous
  expansion or demand-driven expansion, charging access computation, and
  adjusted access to congestion. In each step time $t$, it is proceeded as
  follows:
  \begin{enumerate}
  \item[(a)] Expand charging infrastructure (exogenous or demand-driven). In
    the mode exogenous a number of new chargers is added by means the
    charger expansion rate. In the mode demand-driven the expected number of
    new chargers is determined by $\lambda_t = r\cdot g\cdot s_{EV,t}$,
    where $r$ is the rollout rate, $g$ a grid-size scaling factor, and
    $s_{EV,t}$ is the EV adoption share at time $t$. It means that spatially
    weighted placement based on adopter density and lack of access.
  \item[(b)] Update charging access and effective charging access (in the
    case of congestion). For each cell $(x,y)$, it is computed the baseline
    charging access by
    \begin{equation}
      A_t(x,y) = \Big(1+\frac{d_t(x,y)}{\delta}\Big)^{-1},
    \end{equation}
    where $d_t(x,y)$ is the toroidal Manhattan distance to the nearest
    charger, and $\delta$ is a decay parameter. Effective charging access
    is updated by $A^{e\!f\!f}_t(x,y) = A_t(x,y)\, F_t(x,y)$, with
    \begin{equation}
      F_t(x,y) = \min\Big(1, \frac{\mathrm{Cap}_t(x,y)}{D_{R,t}(x,y)}\Big),
    \end{equation}
    where $D_{R,t}(x,y)$ is the local EV demand (within a neighborhood with
    radius $R$), and $\mathrm{Cap}_t(x,y)$ is the local capacity related
    with local charger supply.
  \end{enumerate}

  Finally, the model records aggregate indicators such as adoption share,
  charging accessibility, and TCO gaps.

\item \textbf{Agent's decision rule.} The decision rule (probabilistic or
  deterministic) is defined by a EV adoption score obtained from its utility
  function integrating economic variables, charging accessibility,
  environmental preferences, peer influence, and range anxiety, combined
  with a hard affordability constraint (income-budget gate,
  Section~\ref{sec:decision}).

  EV adoption score is composed by five components:
  \begin{itemize}
  \item Economic - Agents compute total cost of ownership (TCO) for EV and
    ICE vehicles using effective purchase price (including subsidies),
    energy cost, maintenance, annual mileage, and evaluation horizon. The
    economic component combines TCO advantage and price sensitivity.
  \item Contextual - Charging infrastructural access is a weighted linear
    combination of home charging availability and public charging access at
    the agent's fixed residential location.
  \item Environmental - Environmental awareness combines the specific
    (regarding EVs) environmental concern with the generic (regarding
    affordances) agent's pro-environmental state.
  \item Social (Peer exposure) - Peer influence is proportional to the
    fraction of visible peers (network or residential) who have adopted
    EVs, modulated by peer sensitivity. Social diffusion may also reduce
    range anxiety and increase environmental concern.
  \item Range Anxiety - Perceived range anxiety is a penalty to EV
    adoption, representing then a psychological barrier that is alleviated
    by good effective charging access and through peer exposure.
  \end{itemize}

  In addition to the score, adoption requires \emph{affordability}: the
  average annual cost of EV ownership may not exceed a fixed share of the
  agent's income (Section~\ref{sec:decision}).

  EV adoption only is evaluated when the agent's vehicle reaches its
  replacement interval, ensuring that adoption is tied to realistic
  replacement cycles rather than continuous choice.
\end{enumerate}

\subsubsection{Socio-economic and behavioural heterogeneity}

\begin{table}[ht]
\centering
\caption{Agents attribute distributions. Income is anchored to empirical
data (see text); remaining behavioural traits are sampled from bounded
intervals.}
\label{tab:attributes}
\begin{tabular}{lll}
\toprule
\textbf{Attribute} & \textbf{Distribution} & \textbf{Description} \\
\midrule
Income ($income_i$) & LogNormal($\mu=11600$, $\sigma=7600$) &
  Annual income (EUR), empirically anchored. \\
Annual mileage ($m_i$) & $N(\mu=12000,\ \sigma=2000)$ &
  Yearly distance travelled (km). \\
Replacement interval ($r_i$) & Uniform$(6, 14)$ &
  Years between vehicle replacements. \\
Vehicle age ($a_i$) & Uniform$(0, r_i - 1)$ &
  Initial age, staggered below own $r_i$. \\
Home charging access ($h_i$) & Uniform$(0, 1)$ &
  Probability or ease of charging at home. \\
Environmental concern ($env_i$) & Uniform$(0, 1)$ &
  Specific environmental concern. \\
Price sensitivity ($s^{price}_i$) & Uniform$(0.5, 1.5)$ &
  Weight placed on economic considerations. \\
Peer sensitivity ($s^{peer}_i$) & Uniform$(0.5, 1.5)$ &
  Susceptibility to peer influence. \\
Range anxiety ($r^{anx}_i$) & Uniform$(0, 1)$ &
  Fear of running out of battery. \\
\bottomrule
\end{tabular}
\end{table}

Agents are heterogeneous including socio-economic attributes such as income,
annual mileage, vehicle age, replacement interval, and five behavioural
attributes: home charging access, environmental concern, price sensitivity,
range anxiety, and peer sensitivity. Agents also have behavioural state
(pro-environmental or non-environmental) inherited from the base model,
enabling interactions between environmental cues and EV adoption. The
agents are initialized through distributions that determine agents behaviour
(see Table~\ref{tab:attributes}).

\paragraph{Empirical income anchor.} The income distribution is not a free
parameter: it reproduces Eurostat EU-SILC mean equivalised net disposable
income for Portugal (dataset \texttt{ilc\_di03}, EUR), whose 2010--2024
average is EUR~11{,}565. The lognormal shape parameter
$\sigma_{\ln}\approx 0.60$ is supported by two independent derivations: the
observed mean/median ratio ($\approx 1.20$) and the Gini coefficient of
equivalised disposable income (\texttt{ilc\_di12}, $\approx 33$ over the
period). The parameterization (mean EUR~11{,}600, standard deviation
EUR~7{,}600) reproduces both the empirical mean and median
($\approx$EUR~9{,}700 vs.\ observed EUR~9{,}656). Equivalised income
is expressed per adult-equivalent; the household equivalence-scale factor is
absorbed by the income-budget share (Section~\ref{sec:decision}), which is
the parameter fitted in calibration instead of income.

Initial vehicle age is staggered below each agent's own replacement interval
so that the population does not evaluate adoption in a synchronized burst at
the first step. Price and peer sensitivities are drawn directly on
$[0.5, 1.5]$ and multiply their score components (equivalent to the earlier
$0.5+U(0,1)$ formulation, with a simpler expression).

In addition to public charging accessibility, agents have availability of
charging at own residence, being representing by parameter $h_i$. Higher
values indicating greater convenience (e.g., private garage with installed
charger) while lower values indicating limited or no home charging
capability (e.g., on-street parking without infrastructure). Agents has
preferences toward environmental issues (e.g., perceived reducing
emissions), represented by parameter $env_i$. Agents are assigned a range
anxiety parameter $r^{anx}_i$, representing the perceived risk of running
out of battery before reaching a charging site. The parameter $s^{peer}_i$
describes how strongly the agent responds to EV adoption behaviour of its
neighbours.

\subsubsection{Environment and spatial structure}

The spatial layers provide the environmental and infrastructural context for
agents behaviour. The affordance landscape shapes environmental attitudes,
while the charging access layer shapes EV adoption feasibility. Their
combination enables spatially heterogeneous diffusion patterns and realistic
interactions between infrastructure and behaviour.

\paragraph{Affordance landscape layer.} A central feature of the base model
is the presence of a spatial layer of environmental affordances representing
opportunities or cues that facilitate pro-environmental behaviour. For each
cell in the spatial grid, it is assigned a binary value, which is equal to 1
in the presence of a pro-environmental affordance or equal to 0 otherwise.
These affordances encode the idea that behaviour is shaped not only by
internal perceptions but also by contextual cues that either enable or
inhibit action. Although the EVAM does not directly use affordances as
behavioural triggers, the affordance landscape remains part of the
environment and contributes to the agent's pro-environmental state, which in
turn influences EV adoption score.

\paragraph{Charging infrastructural access layer.} A second spatial layer
represents public charging accessibility. Each cell stores a value in the
interval $[0,1]$ indicating the relative ease of accessing public charging
infrastructure at that location. This layer evolves dynamically over time
through an exogenous expansion process governed by the parameter
\texttt{charger\_expansion\_rate}. The charging access field captures
spatial inequalities in infrastructure availability and provides a
continuous measure of charging convenience that agents incorporate into
their adoption decisions.

\paragraph{Spatial neighbourhood structure.} Peer (spatial or social)
effects reuse the existing network implementation from the base model.
Agents interact with their local environment through a Moore neighbourhood,
which includes eight surrounding cells. When diffusion social option is
enabled, peers are agent's social network neighbours. Otherwise peers are
residential neighbours. Let $N_i$ denotes the set of neighbours of agent
$i$, and $A_j$ be an indicator equal to 1 if neighbour $j$ has already
adopted an EV. Then peer adoption share is given by
\begin{equation}
P_{i,t} =
\begin{cases}
0, & |N_i| = 0,\\[2pt]
\dfrac{1}{|N_i|}\sum_{j\in N_i} A_j, & |N_i| > 0.
\end{cases}
\end{equation}

Peer adoption visibility has two effects:
\begin{enumerate}
\item \textbf{Range anxiety reduction.} Observing EV use among peers reduces
  perceived range anxiety. The reduction is proportional to peer adoption
  share. This reflects empirical evidences that social exposure alleviates
  concerns about driving range.
\item \textbf{Environmental concern increase.} Peer adoption increases the
  salience of environmental motivations specific to EVs. The increase is
  proportional to peer adoption share. This corresponds to social
  reinforcement of pro-environmental norms.
\end{enumerate}
Therefore it represents a deeper form of social learning acting on
psychological traits within the agent population.

\paragraph{Agent placement and mobility.} Agents are initially placed at
random locations on the grid. Since then agents have a fixed residential
location and move to start pro-environmental learning dynamics driven by
affordances. It is highlighted that EV adoption is strongly tied to
residential location, which determines access to home and public charging
infrastructure and local peer exposure.

\subsubsection{Agent's decision-making process}
\label{sec:decision}

Agent's decision-making process is based on the decision rule (probabilistic
or deterministic) which depends of a EV adoption score, combined with a hard
affordability constraint.

Agent's EV adoption score is a multi-dimensional metric that drives EV
adoption dynamics. For each agent $i$ at time $t$, it is defined by
\begin{equation}
Score_{i,t} = w_{econ} Econ_{i,t} + w_{charge} C_{i,t} + w_{env} Env_{i,t}
            + w_{peer} Peer_{i,t} - w_{range} R_{i,t}
\end{equation}
where $w_{econ}$, $w_{charge}$, $w_{env}$, $w_{peer}$, and $w_{range}$ are
parameters.

The deterministic decision rule is represented by
\begin{equation}
d_{i,t} =
\begin{cases}
1, & \text{if } Score_{i,t} \geq \theta,\\
0, & \text{otherwise}
\end{cases}
\end{equation}
in which $\theta$ is a threshold parameter. In turn, the probabilistic
decision rule is given by
\begin{equation}
p_{i,t} = \frac{1}{1+e^{-(Score_{i,t}-\theta)/T}},
\end{equation}
where the parameter $T$ controls the steepness of the logistic curve. Higher
$T$ values correspond to greater behavioural uncertainty.

\paragraph{Income-budget gate (affordability).} Independently of the
score-based decision, an agent may only adopt if the average annual cost of
EV ownership does not exceed a share $\beta$ of its income:
\begin{equation}
\frac{TCO_{EV,i,t}}{Y} \;\leq\; \beta \cdot income_i .
\label{eq:budgetgate}
\end{equation}
This implements the assumption that households cannot devote more than a
fixed fraction of annual income to vehicle ownership as an actual
constraint. With income pinned to the empirical Eurostat distribution, the
calibration (Section~\ref{sec:results}) fitted $\beta = 0.108$ ---
essentially recovering the familiar ``10\% of income'' budget heuristic
from the adoption data, and absorbing the household equivalence-scale
factor of the income measure.

We now define the referred components (without weights) of the agent's EV
adoption score.

\paragraph{Economic component.} The economic component at time $t$ is
calculated by
\begin{equation}
Econ_{i,t} = s^{price}_i \, \frac{\Delta(TCO_{i,t})}{TCO_{ICE,i}},
\end{equation}
with $s^{price}_i \sim U(0.5, 1.5)$. Affordability is not part of the
score; it enters as the hard gate of Eq.~\eqref{eq:budgetgate}.

\paragraph{Charging infrastructural access component.} For each agent $i$
located at home position $(x,y)$ and time $t$,
\begin{equation}
C_{i,t}(x,y) = 0.7\, h_i + 0.3\, A^{e\!f\!f}_{i,t}(x,y)
\end{equation}
where $h_i$ denotes home charging access and $A^{e\!f\!f}$ the effective
charging access.

\paragraph{Environmental component.} For each agent $i$ at time $t$,
\begin{equation}
Env_{i,t} = (1-w)\, env_{i,t} + w\, proenv_{i,t}
\end{equation}
where $w$ is a weight parameter, $proenv_{i,t}$ is the generic
pro-environmental behavioural state, and $env_{i,t}$ is the specific
environmental concern. The latter remains fixed unless social diffusion is
enabled, in which case peer exposure may update.

\paragraph{Social (peer exposure) component.} For each agent $i$ at time
$t$,
\begin{equation}
Peer_{i,t} = P_{i,t}\; s^{peer}_i ,
\end{equation}
where $s^{peer}_i \sim U(0.5, 1.5)$ is peer sensitivity and $P_{i,t}$ is
peer adoption share.

\paragraph{Range anxiety component (penalty).} Range anxiety component is
given by
\begin{equation}
R_{i,t} = s^{range}_{i,t}\big(1 - A^{e\!f\!f}_{i,t}(x,y)\big),
\end{equation}
where $s^{range}_{i,t}$ is range anxiety and $A^{e\!f\!f}$ denotes the
effective charging access. Observe that range anxiety penalises the EV
adoption if effective charging access is low.

\section{Simulation results}
\label{sec:results}

Model parameters are defined through the \texttt{EVParams} class, which
provides baseline values for economic variables (fuel price, EV purchase
cost, ICE operating cost), infrastructure variables (charger expansion rate,
initial charging coverage), and behavioral coefficients. Scenario-specific
overrides are applied using the \texttt{EVParams.from\_scenario} method,
allowing controlled variation of selected parameters while keeping all
other model components fixed.

\subsection{Calibration to Portugal (2010--2020) and hold-out validation
(2021--2024)}
\label{sec:calibration}

The model is calibrated against the Portuguese battery-electric vehicle
(BEV) \emph{fleet} share --- the fraction of the circulating light-vehicle
fleet that is electric, constructed from UVE/IMT fleet counts over a
light-vehicle fleet of $\approx$5.8 million (note this is roughly an order
of magnitude smaller than the commonly reported \emph{new-registration}
share). The scenario spans annual steps 2010--2024, applies the historical
subsidy trajectory as a schedule (zero through 2014, then a ramp from
EUR~500 to EUR~4{,}000 reflecting the 2015 reintroduction of purchase
incentives and the documented grant range), and seeds the observed 2010
stock.

Calibration deliberately uses only 2010--2020 (log-space RMSE objective,
weighting the early exponential phase that raw RMSE ignores);
\textbf{2021--2024 is a hold-out never seen by the search}. With income
pinned to the empirical Eurostat distribution, the search fits the
income-budget share $\beta$, vehicle prices, the Wright learning rate, the
charger rollout rate, and the score threshold. Figure~\ref{fig:calibration}
shows the result over 12 seeds: fit-window log-RMSE 0.36 (raw RMSE 0.0006);
hold-out log-RMSE 0.29 (raw RMSE 0.006). Hold-out error is comparable to
in-sample error, i.e.\ there is no evidence of overfitting.

\begin{figure}[ht]
\centering
\includegraphics[width=0.85\textwidth]{outputs/v13_figures/fig_calibration.png}
\caption{Portugal BEV fleet share, model (12-seed mean $\pm$1 sd) versus
observed. The model is fitted on 2010--2020 only; 2021--2024 is out of
sample.}
\label{fig:calibration}
\end{figure}

Two structural findings emerge. First, the calibrated regime is
\emph{budget-gate dominated}: the fitted score threshold is nearly always
passed, and adoption timing is governed by the affordability constraint
interacting with the subsidy ramp and the Wright price decline, together
with charging access. Second, out of sample the model systematically
\emph{undershoots the post-2020 surge} (2024: model $\approx$2.3\% vs.\
observed 3.28\%): the drivers of that surge --- the explosion of available
EV models, corporate fleet electrification, and the 2022 fuel-price spike
--- are outside the model's scope. We report this as a scope statement
rather than a fitting failure.

\subsection{Policy counterfactuals on the calibrated scenario}

Rather than abstract scenario presets, policy experiments are run as
counterfactuals on the calibrated Portugal scenario (4{,}000 agents, 12
seeds, mean $\pm$1 sd), asking how the 2024 fleet share would have differed
under alternative policies (Figure~\ref{fig:scenarios}).

\begin{figure}[ht]
\centering
\includegraphics[width=0.9\textwidth]{outputs/v13_figures/fig_scenarios.png}
\caption{BEV fleet share reached in 2024 under policy counterfactuals
(12 seeds, error bars $\pm$1 sd).}
\label{fig:scenarios}
\end{figure}

The subsidy trajectory is the decisive lever. Removing the subsidy ramp
collapses the 2024 fleet share from $2.3\% \pm 1.6$ to $0.3\% \pm 0.6$ ---
without purchase incentives the affordability gate keeps essentially the
entire income distribution out of the market and the Wright price decline
never ignites. Doubling the ramp (peak EUR~8{,}000) nearly triples the 2024
share to $6.4\% \pm 0.9$, and also stabilises the trajectory across seeds
(the dispersion halves), because a deeper subsidy reaches further down the
steep part of the income distribution and starts the price-learning
feedback earlier and more reliably. By contrast, doubling the charger
rollout leaves the outcome unchanged ($2.3\% \pm 1.6$, identical to
baseline), and a high fuel price (EUR~2.6/l) adds only $\approx$0.1
percentage points: in the calibrated, budget-gate-dominated regime,
neither infrastructure nor operating-cost levers relax the constraint that
actually binds.

\subsection{Sensitivity analysis}

All sensitivity sweeps are run at the calibrated scale with 8 seeds per
point and report mean $\pm$1 sd of the 2024 fleet share; single-run curves
(as in earlier drafts) proved unstable across seeds and are not
interpreted.

\begin{figure}[ht]
\centering
\includegraphics[width=0.75\textwidth]{outputs/v13_figures/fig_sens_fuel.png}
\caption{2024 BEV fleet share as a function of fuel price.}
\label{fig:fuel}
\end{figure}

\begin{figure}[ht]
\centering
\includegraphics[width=0.75\textwidth]{outputs/v13_figures/fig_sens_charger.png}
\caption{2024 BEV fleet share as a function of charger expansion rate.}
\label{fig:charger}
\end{figure}

\begin{figure}[ht]
\centering
\includegraphics[width=0.75\textwidth]{outputs/v13_figures/fig_sens_subsidy.png}
\caption{2024 BEV fleet share as a function of the subsidy-ramp scale
(x-axis: peak subsidy reached by 2024).}
\label{fig:subsidy}
\end{figure}

\begin{figure}[ht]
\centering
\includegraphics[width=0.85\textwidth]{outputs/v13_figures/fig_surface.png}
\caption{2024 BEV fleet share as a function of fuel price and charger
expansion rate (4-seed means).}
\label{fig:surface}
\end{figure}

The sweeps confirm the counterfactual picture. The subsidy response
(Figure~\ref{fig:subsidy}) is strong and convex --- from $0.4\%$ with no
subsidy to $6.3\%$ at a doubled ramp --- because each additional
EUR~1{,}000 of subsidy simultaneously lowers the affordability cutoff into
denser regions of the income distribution and accelerates the Wright
price decline. The fuel-price response (Figure~\ref{fig:fuel}) is positive
but modest ($1.4\%$ at EUR~1.4/l to $2.3\%$ at EUR~3.0/l): operating costs
shift the TCO comparison but not the binding budget constraint, which
depends primarily on the purchase price net of subsidy. The
charger-expansion response (Figure~\ref{fig:charger}) is nearly flat
($1.6\%$ at zero rollout to $2.0\%$ at rate 6), and the joint surface
(Figure~\ref{fig:surface}) is correspondingly gentle (from $1.1\%$ at the
low-low corner to $2.6\%$ at the high-high corner) with no interaction
structure that survives seed averaging. We emphasise what is \emph{absent}
relative to earlier drafts: with 8--12 seeds per point, none of the
apparent threshold effects, tipping zones, or non-monotonicities visible
in single-run curves survive; those features were seed noise.

\subsection{Policy implications}

The calibrated model carries three implications for Portuguese-style EV
policy. First, purchase-price support dominates: in a market where the
binding constraint is the household budget, subsidies act twice --- once
directly on affordability and once indirectly by igniting the
learning-curve price decline --- and their effect is convex over the
historically relevant range. Second, charging infrastructure and fuel
prices, prominent levers in the earlier exploratory analysis, are
second-order in the calibrated regime: they shape the adoption score, but
the score threshold is rarely the binding margin. This does not imply
infrastructure is unimportant in reality --- rather, that over Portugal's
2010--2024 fleet-share trajectory, affordability was the bottleneck the
data identify. Third, the out-of-sample undershoot of the 2021--2024 surge
(Section~\ref{sec:calibration}) marks the model's scope boundary: drivers
such as the expanding EV model range, corporate fleet electrification, and
the 2022 fuel-price shock are not represented, and extending the model to
capture them is the natural next step.

\appendix
\section*{Appendix A: ODD Protocol for the EV Adoption Model (EVAM)}

\subsection*{A.1 Overview}

Our model description follows the ODD (Overview, Design concepts, Details)
protocol for describing agent-based models
(\cite{grimm2006,grimm2010,grimm2020}).

\begin{enumerate}
\item \textbf{Purpose and patterns.} The Electric Vehicle Adoption Model
  (EVAM), implemented in a Python/Mesa framework, extends the original
  NetLogo Affordance Landscape Model (\cite{kaaronen2020}) by adding
  household EV adoption. EVAM is designed to explore the behavioural,
  infrastructural, and economic mechanisms that shape the adoption of
  electric vehicles. The adoption emerges from total cost of ownership
  (TCO), charging infrastructure availability (home and public charging
  access), environmental attitudes, social influence (peer effects), range
  anxiety, and a hard affordability constraint, under policy levers such as
  subsidies, fuel prices, charging roll-out, and optional market feedbacks
  (demand-driven infrastructure, congestion, supply limits, price learning,
  and social diffusion). The patterns the model is expected to reproduce
  are: (i) the near-zero early phase and subsequent exponential growth of
  the Portuguese BEV fleet share (2010--2020, used for calibration); and
  (ii) plausible out-of-sample continuation (2021--2024, used for hold-out
  validation). It is intended for scenario exploration, policy analysis,
  and the study of spatial diffusion patterns.

\item \textbf{Entities, state variables, and scales.} There are three types
  of agents in the EVAM: households, represented by mobile circle-shaped
  agents, affordances (static patches that occupy grid cells) and links
  (which connect agents in a social network).

  \emph{Space.} The environment is a toroidal Mesa Multigrid of size
  $Width\times Height$ shared with the base model, with three relevant
  layers: the affordance landscape (unchanged from the base model); a
  charging access layer $A_t(x,y)=(1+d_t(x,y)/\delta)^{-1}$ where $d$ is
  the Manhattan distance to the nearest charger site; and an effective
  charging access layer $A^{e\!f\!f}_t = A_t F_t$ with congestion factor
  $F_t=\min(1,\mathrm{Cap}_t/D_{R,t})$, identical to charging access when
  charger capacity is infinite. Charger sites are point locations with at
  most one charger per cell.

  \emph{Agents.} EVAM includes an agent class (\texttt{EVConsumerAgent}), a
  subclass of \texttt{ConsumerAgent} of the base model. Agents are
  characterised by socio-economic attributes (income, annual mileage,
  vehicle age, vehicle replacement interval); behavioural traits
  (environmental concern, price sensitivity, peer sensitivity, range
  anxiety); and contextual attributes (home and public charging access).
  Dynamic affordance states (\texttt{pro\_env}, \texttt{non\_env}) are
  inherited from the base model and are updated by affordance learning.
  Agents have two positions with distinct meanings: \texttt{pos} --- the
  moving position in the base model's behaviour loop (not residential
  relocation); and \texttt{home\_pos} --- the fixed residential location
  assigned at placement, used by all EV mechanics (charging access reads,
  spatial peer effects, charger demand). The adoption state
  (\texttt{ev\_adopted}) is one-shot and adopting resets vehicle age to
  zero. Time advances in discrete yearly steps. Agents do not relocate.

  EVAM records decision telemetry on every evaluation:
  \texttt{last\_adoption\_score}, \texttt{last\_economic\_score},
  \texttt{last\_charging\_score}, \texttt{last\_environmental\_score},
  \texttt{last\_peer\_adoption\_share},
  \texttt{last\_range\_anxiety\_penalty}, \texttt{last\_ice\_tco},
  \texttt{last\_tco\_gap}, \texttt{last\_ev\_tco},
  \texttt{last\_adoption\_probability}, \texttt{last\_affordable} (the
  income-budget gate outcome), and \texttt{has\_evaluated\_adoption}.

  \emph{Model-level state.} Beyond agent-level states, EVAM contains global
  state variables related to EV, including charger locations, the effective
  EV price (learning curve), the per-step resolved policy/price values
  (\texttt{current\_subsidy}, \texttt{current\_fuel\_price},
  \texttt{current\_electricity\_price}), per-step market counters, and
  aggregated metrics.

\item \textbf{Process overview and scheduling.} Each simulation step:
  (a) reset market counters, set the effective price from the previous
  step's cumulative adoption, and resolve the current subsidy and energy
  prices from their optional schedules; (b) expand charging infrastructure
  (exogenous or demand-driven); (c) update charging access and effective
  charging access (congestion); (d) activate agents in random order to
  perform the affordance behaviour loop (move/act/learn), age vehicle,
  optionally diffuse perceptions, and evaluate EV adoption if the
  replacement interval is reached (score rule + income-budget gate + market
  supply gate); (e) update EV metrics and record the full model state. Row
  0 of the DataCollector is collected at construction.
\end{enumerate}

\subsection*{A.2 Design concepts}

\begin{description}
\item[Basic principles.] EVAM integrates ecological psychology
  (affordances), behavioural economics (TCO, price sensitivity, budget
  constraint), and social diffusion. Adoption emerges from
  multi-dimensional behavioural drivers modulated by infrastructure and
  policy.
\item[Emergence.] EVAM produces spatial clustering of EV adoption,
  diffusion waves driven by infrastructure rollout, congestion effects
  reducing perceived access, socio-spatial inequalities, and peer-driven
  acceleration or stagnation of adoption.
\item[Adaptation.] Agents adapt annually by updating their adoption score
  based on current economic, infrastructural, environmental, and social
  conditions.
\item[Objectives.] Agents aim to minimise long-term vehicle ownership costs
  while satisfying behavioural preferences, social norms, and their budget
  constraint.
\item[Learning.] Agents do not learn over time in the reinforcement sense.
  However, environmental concern and range anxiety may shift
  deterministically through social diffusion, hence peer influence acts as
  a form of social learning.
\item[Prediction.] Agents do not forecast future conditions so that
  decisions rely on current-year information.
\item[Sensing.] Agents observe the economic market (TCO, current fuel and
  electricity price), policy incentives (current subsidy), public charging
  access at their home location, availability of home charging, peer
  adoption share, their own traits, and affordance states.
\item[Interaction.] Agents interact through peer influence, either via
  social networks or residential neighbourhoods (Moore-1 torus).
\item[Stochasticity.] Initial trait sampling (including lognormal income),
  initial adopter assignment, random activation order, random charger
  placement in exogenous mode, and the logistic decision rule when enabled.
\item[Collectives.] Neighbourhoods (spatial or network-based) act as
  collectives influencing adoption through peer visibility.
\item[Observation.] EVAM records EV adoption share, charger count,
  effective EV price, current policy/price values, congestion, TCO gap,
  component scores, affordability outcomes, and agent-level telemetry.
\end{description}

\subsection*{A.3 Details}

\begin{enumerate}
\setcounter{enumi}{4}
\item \textbf{Initialization.} Agents are placed randomly on the grid and
  endowed with socio-economic attributes and behavioural traits
  (Table~\ref{tab:attributes}); initial vehicle age is staggered below each
  agent's replacement interval. Initial adopters are assigned according to
  parameter-defined share (the Portugal scenario seeds the observed 2010
  stock). Initial chargers are placed, charging access and effective access
  are computed, the effective EV price is set by the learning curve, and
  the DataCollector records the initial state.

\item \textbf{Input data.} The income distribution is anchored to Eurostat
  EU-SILC statistics for Portugal (\texttt{ilc\_di03} mean/median
  equivalised net disposable income; \texttt{ilc\_di12} Gini coefficient).
  The calibration target is the Portuguese BEV fleet-share series
  2010--2024 constructed from UVE/IMT fleet counts and an ACAP/APA fleet
  denominator, and the Portugal scenario applies the historical subsidy
  trajectory as an exogenous schedule. All remaining parameters are
  user-defined through \texttt{EVParams} or scenario presets.

\item \textbf{Submodels.} The agent's EV adoption score aggregates
  economic, infrastructural, environmental, social, and behavioural
  components as defined in Section~\ref{sec:decision}, using the Wright's
  law effective price, the (optionally discounted) TCO of
  Section~3.2, the income-budget gate of Eq.~\eqref{eq:budgetgate}, and
  the market supply constraint
  $\sum_i \mathbf{1}\{\mathrm{adopt}_i(t)\} \leq
  \texttt{ev\_supply\_per\_step}$.
\end{enumerate}

\paragraph{Reproducibility.} All stochasticity flows through model-seeded
RNGs. Disabling optional mechanisms yields deterministic, byte-identical
results for fixed seeds. All simulation scripts, parameter files, and
notebooks used in this analysis are stored in the project repository
(including \texttt{scripts/calibrate\_portugal.py} for the calibration
search and \texttt{scripts/make\_v13\_figures.py}, which regenerates every
figure in Section~\ref{sec:results}), ensuring reproducibility of the
results.

\begin{thebibliography}{12}
\bibitem{bonabeau2002} Bonabeau, E. (2002). Agent-based modeling: Methods
and techniques for simulating human systems, \emph{Proc. Natl. Acad. Sci.
U.S.A.} 99 (3) 7280--7287.
\url{https://doi.org/10.1073/pnas.082080899}
\bibitem{elkington1997} Elkington, J. (1997). \emph{Cannibals with Forks:
The Triple Bottom Line of 21st Century Business}. Capstone, Oxford.
\bibitem{grimm2006} Grimm, V., Berger, U., Bastiansen, F., et al. (2006). A
standard protocol for describing individual-based and agent-based models.
\emph{Ecological Modelling}, 198(1--2), 115--126.
\url{https://doi.org/10.1016/j.ecolmodel.2006.04.023}
\bibitem{grimm2010} Grimm, V., Berger, U., DeAngelis, D.~L., Polhill,
J.~G., Giske, J., \& Railsback, S.~F. (2010). The ODD protocol: A review
and first update. \emph{Ecological Modelling}, 221(23), 2760--2768.
\url{https://doi.org/10.1016/j.ecolmodel.2010.08.019}
\bibitem{grimm2020} Grimm, V., Railsback, S.~F., Vincenot, C.~E., et al.
(2020). The ODD protocol for describing agent-based and other simulation
models: A second update to improve clarity, replication, and structural
realism. \emph{J. Artif. Soc. Soc. Simul.} 23(2).
\url{http://jasss.soc.surrey.ac.uk/23/2/7.html}
\bibitem{kaaronen2020} Kaaronen, R.~O., \& Strelkovskii, N. (2020).
Cultural Evolution of Sustainable Behaviors: Pro-environmental Tipping
Points in an Agent-Based Model. \emph{One Earth}, 2(1), 85--97.
\url{https://doi.org/10.1016/j.oneear.2020.01.003}
\bibitem{mohammed2024} Mohammed, A., Saif, O., Abo-Adma, M., et al. (2024).
Strategies and sustainability in fast charging station deployment for
electric vehicles. \emph{Scientific Reports} 14, 283.
\url{https://doi.org/10.1038/s41598-023-50825-7}
\bibitem{nielsen2024} Nielsen, K.~S., Cologna, V., Bauer, J.~M., et al.
(2024). Realizing the full potential of behavioural science for climate
change mitigation. \emph{Nat. Clim. Chang.} 14, 322--330.
\bibitem{rai2016} Rai, V., \& Henry, A.~D. (2016). Agent-based modelling of
consumer energy choices. \emph{Nat. Clim. Change} 6, 556--562.
\bibitem{orourke2015} O'Rourke, D., \& Lollo, N. (2015). Transforming
Consumption: From Decoupling, to Behavior Change, to System Changes for
Sustainable Consumption. \emph{Annual Review of Environment and Resources},
40, 233--259. \url{https://doi.org/10.1146/annurev-environ-102014-021224}
\bibitem{syed2024} Syed, S., Acquaye, A., Khalfan, M.~M., Obuobisa-Darko,
T., \& Yamoah, F. (2024). Decoding Sustainable Consumption Behavior: A
Systematic Review of Theories and Models and Provision of a Guidance
Framework. \emph{Resources, Conservation \& Recycling Advances}.
\bibitem{eurostat} Eurostat (2026). EU-SILC: mean and median equivalised
net disposable income (\texttt{ilc\_di03}); Gini coefficient of equivalised
disposable income (\texttt{ilc\_di12}). Retrieved July 2026.
\bibitem{uve} UVE --- Utilizadores de Ve\'iculos El\'etricos (2025).
Parque de Ve\'iculos El\'etricos em Portugal.
\url{https://www.uve.pt/page/parque-ve-2024/}
\end{thebibliography}

\end{document}
